\section{Organização inicial dos dados}
\subsection{Requisitos para assuntos de informação}
\begin{table}[!ht]\centering\small\renewcommand{\arraystretch}{1.45}
\begin{tabularx}{\linewidth}{|>{\raggedright\arraybackslash}X|>{\raggedright\arraybackslash}X|>{\centering\arraybackslash}p{2cm}|}
\hline
\textbf{Assunto / grupo} & \textbf{Características} & \textbf{RFs} \\
\hline
Aviso & ID, título, conteúdo, categoria, status, criação, publicação & RF01--RF06 \\
\hline
Usuário & ID, nome, papel, credenciais, nível de acesso & RF01, RF03--RF06 \\
\hline
Categoria de Aviso & ID, nome, indicador de moderação & RF02 \\
\hline
Solicitação de autorização & ID, aviso, avaliador, decisão, feedback, datas & RF03--RF05 \\
\hline
Avaliação & ID, parecer, decisão, data e status & RF04--RF05 \\
\hline
Público-alvo & ID, descrição, turma, função ou seleção específica & RF06 \\
\hline
Envio & Aviso, canal, data, tipo de segmentação & RF03--RF06 \\
\hline
\end{tabularx}
\end{table}

\subsection{Conexões entre os assuntos}
\begin{itemize}
\item Um emissor cria um ou mais avisos.
\item Um aviso pode possuir uma ou mais solicitações de autorização.
\item Um avaliador realiza uma avaliação para um aviso pendente.
\item Um aviso é direcionado a um ou mais públicos-alvo.
\item Um destinatário visualiza avisos vinculados ao seu público.
\end{itemize}

\subsection{Informações ainda a confirmar}
\begin{enumerate}
\item Quais categorias exigem moderação obrigatória?
\item Como os vínculos entre responsáveis, alunos e turmas devem ser cadastrados?
\item É necessário confirmar digitalmente a leitura de avisos críticos?
\item Quais atributos e estados do processo precisam ser persistidos para auditoria?
\end{enumerate}
